\chapter{Closing the cavity inequality}\label{sec:c7-closure}

In this section, we complete the cavity argument of
\cref{sec:c7,sec:c7-endpoint}. The main result is
\cref{prop:c7-closure}: if the scalar response operator at a minimizing
profile satisfies $J^{\scal}_{a,d}\preceq(1+\eps_N)I$ with
$0\le\eps_N\le\rho/4$, then the self-overlap variance satisfies
$(3\rho/4)V\le C_{\rm cav}[N^{-1}+A+J+H_3]$. In \cref{cor:c7-defect},
we absorb the cubic moment $H_3$ and bound the defect $\calD_N(\rho)$ by
$C_\beta\log^2(2+N)/(N\lambda\rho^2)$ for $1/N\le\rho\le1$, wherever
the defect is defined and some minimizer satisfies the stability
assumption. In \cref{sec:b7}, we use this corollary to bound the contribution of
the stable interval in the proof of \cref{thm:main}, and we verify the
stability assumption there with \cref{prop:lz-stability}.

We rely on the following earlier results. \Cref{sec:c7}
provides the laws of finitely many descendants with
\cref{lem:streaming}, the comparison of nonnegative observables in
\cref{cor:comparison}, the three-defect bound of
\cref{lem:cav-threeproducts}, and the Taylor remainder of
\cref{lem:remainder}. \Cref{sec:c7-endpoint} provides the integrated
linear response of \cref{lem:linear} and the signed response estimate
of \cref{lem:cl-response}, and \cref{lem:Jdom} states that
$J^{\scal}_{a,d}$ is bounded and self-adjoint. We also use the contact
identities
\eqref{eq:in-contact}, the conditional decomposition in the proof of
\cref{lem:fnd-decomp}, the trace bound \eqref{eq:in-Lvar}, the bound on
$J$ in \cref{prop:th-j}, the cubic moment bound of \cref{cor:wi-cub},
and the defect identity \eqref{eq:in-envelope}.

\paragraph{Standing assumptions and notation.} We keep the setting of
\cref{sec:c7}: $N\ge2$, $n=N-1$, $0<\rho\le1$, $t=1-\rho$,
$c=\beta^2\rho$, and $0<s=\lambda/\beta\le1/2$. The parameters $\alpha$, $d$, and
$b(u)$ are given by \eqref{eq:cavityparameters}. The profile
$\mu=s\delta_0+\nu$ is deterministic, with cumulative function $a$, and
each statement says whether $\mu$ must minimize $\mathcal F_{t,s}$ or
satisfy $\nu(\{1\})=0$. Constants $C_\beta$ depend only on $\beta$; they
do not depend on $N$, $\rho$, $s$, the profile, or the positions and
masses of its atoms.

We write $\E_1$, $\Var_1$, and $\Cov_1$ for expectation, variance, and
covariance under $\Pp_u$ at $u=1$. Recall from \cref{subsec:c7-row}
that $\Pp_1$ is the tilted law of \cref{subsec:fnd-law} for the full
system, so that $\E_1=\E$. At $u=0$, we write $\E^-$ and $\Var^-$ for
the rest system, as in \cref{sec:c7-endpoint}. We also recall the rest
self-overlap $A^-_p=|M_p|^2/n$, the conditional magnetization
$m_N(p)$ of the last spin, the descendant laws $\E_{k,q,u}$, the overlap
defects $\Delta_{\mathsf{ab}}$ and $D_{\mathsf{ab}}$ of
\eqref{eq:defectdefinition}, and the integrated third moment
$\mathcal M_3(u)$ of \eqref{eq:cav-M3}. In \cref{subsec:c7-three}, we
defined the covariance
$\mathcal C_{pp}(u)=\Cov_u(m_N(p)^2,A^-_p)$ and its Taylor remainder
$\mathcal R_{\rm cav}$. In \cref{subsec:c7-linear}, we defined the
covariance $\Sigma^-$ of the rest self-overlap, the centered rest
self-overlap $\bar A^-_p$, the error terms $T_p$ of
\eqref{eq:cav-T1}--\eqref{eq:cav-T3}, and the quantity $\mathcal P$ of
\eqref{eq:Ppositive}. \Cref{lem:linear} uses the rest variance
$V^-=\int\Var^-(A^-_p)\,\nu(\dd p)$ and shows that $\mathcal P\ge0$.

\paragraph{Exchangeability.} The tilted law of \cref{subsec:fnd-law} is
invariant under a simultaneous permutation of the $N$ coordinates of
the couplings, the auxiliary field, and the terminal spin. Indeed, let
$\pi$ be a permutation of $\{1,\dots,N\}$, and let $W^\pi$ and $x^\pi$
denote the matrix and the vector with permuted coordinates. The
reference law of $W$ is invariant under $W\mapsto W^\pi$, and since the
terminal condition in \eqref{eq:fnd-vector-pde} and the Laplacian are
invariant, uniqueness gives $U(q,x^\pi;W^\pi)=U(q,x;W)$. The tilt, the
drift in \eqref{eq:fnd-diffusion}, and the Gibbs law of the terminal
spin transform accordingly. In particular, the vector $(m_1(p),\dots,m_N(p))$ of the
conditional magnetizations $m_i(p)=\E[\sigma_i\mid\calG_p]$ is
exchangeable. As $S_p=N^{-1}\sum_im_i(p)^2$ is a symmetric function of
this vector, the pairs $(S_p,m_i(p)^2)$, $1\le i\le N$, all have the
same law. For every deterministic profile and every $p\in[0,1]$, this
implies that
\begin{equation}\label{eq:c7-exchange}
 \E_1m_N(p)^2=\E S_p,\qquad
 \Cov_1\bigl(S_p,m_N(p)^2\bigr)
 =\frac1N\sum_{i=1}^N\Cov\bigl(S_p,m_i(p)^2\bigr)
 =\Var(S_p).
\end{equation}
At $u=1$, the vector $M_p$ consists of the first $n$ coordinates of
$m_p$, and we have
\begin{equation}\label{eq:c7-site-split}
 NS_p=m_N(p)^2+nA^-_p .
\end{equation}
Exchangeability of the sites is used in the same way in the cavity
method of \cite[Section~1.6]{TalagrandVolI2011}.

\section{The mechanism of the contraction}\label{sec:c7-mechanism}

Before giving the proofs, we explain how the inequality closes. In
\cref{lem:cav-trace} below, we prove the identity
\[
 V=\frac{V^{\rm site}}N+\frac nN\Bigl[
 \frac\alpha d\Tr_\nu(J^{\scal}_{a,d}\Sigma^-)
 -\frac\alpha2\mathcal P+\alpha\int T_p\,\nu(\dd p)
 +\mathcal R_{\rm cav}\Bigr],
\]
where $0\le V^{\rm site}\le1/4$. The rest variance $V^-=\Tr\Sigma^-$
enters this identity through the trace pairing
$\Tr_\nu(J^{\scal}_{a,d}\Sigma^-)$. The other terms in the
bracket are either nonpositive, as $-\alpha\mathcal P/2$ is, or errors
controlled by $A$, $J$, $H_3$, and $N^{-1}$
(\cref{lem:cl-response,lem:remainder,lem:c7-thermal}). Under the
stability assumption $J^{\scal}_{a,d}\preceq(1+\eps_N)I$, the first term in
the bracket
is at most $(\alpha/d)(1+\eps_N)V^-$.

The contraction comes from the factor $\alpha/d$. We consider two
descendants $\mathsf a,\mathsf b$ that split at the level
$q_{\mathsf{ab}}$, and fix their rest spins
$\tau^{\mathsf a},\tau^{\mathsf b}$. Under the reference Gaussian law,
the fields acting on $\epsilon_{\mathsf a}$ and $\epsilon_{\mathsf b}$
are the incident row plus the last coordinate of the auxiliary field.
The row contributes covariance $(\alpha u/n)\,\tau^{\mathsf a}\cdot
\tau^{\mathsf b}=\alpha uR^-_{\mathsf{ab}}$, and the last field
contributes $b(u)q_{\mathsf{ab}}$. The total covariance is
\[
 b(u)q_{\mathsf{ab}}+\alpha uR^-_{\mathsf{ab}}
 =d\,q_{\mathsf{ab}}+\alpha u\,(R^-_{\mathsf{ab}}-q_{\mathsf{ab}}).
\]
Only the part $\alpha u(R^-_{\mathsf{ab}}-q_{\mathsf{ab}})$ depends on
the fluctuations of the rest overlap. Its coefficient $\alpha u$ is at
most $\alpha$, while the coefficient of $q_{\mathsf{ab}}$ is $d$. By
\eqref{eq:cavityparameters}, the ratio of the two satisfies
\begin{equation}\label{eq:c7-alpha-d}
 \frac\alpha d=\frac{\beta^2tn/N}{\beta^2(N-t)/N}
 =t\,\frac{N-1}{N-t}\le t=1-\rho .
\end{equation}
The remaining part of the field has deterministic covariance; at
$u=1$, the last field has variance coefficient $b(1)=c=\beta^2\rho$.
Inserting this bound on the first term into the identity above, we
find $V\le t(1+\eps_N)V^-+{}$errors.

To close the inequality, we must compare $V^-$ with $V$, and the
comparison must have coefficient exactly one. Indeed, suppose that we
only knew $V^-\le KV+{}$errors for a constant $K>1$, which is the kind
of bound that a comparison of the two endpoints through
\cref{cor:comparison} gives, with $K=K_{\beta,k}$. We would then have
$V\le tK(1+\eps_N)V+{}$errors, which bounds $V$ only if
$tK(1+\eps_N)<1$. This fails whenever $\rho\le1-1/K$. With coefficient
one, \cref{lem:cl-star} implies that $V^-\le V+{}$errors, and hence
$V\le(1-\rho)(1+\eps_N)V+{}$errors. As
$(1-\rho)(1+\rho/4)\le1-3\rho/4$, a margin of order $\rho$ remains once
$\eps_N\le\rho/4$.

The comparison with coefficient one has two parts. At the independent
endpoint $u=0$, the rest variance is at most the second moment
$\E^-(A^-_p-p)^2$, which is the mean of a product of two overlap defects
on four descendants that split at $p$. The $u$-derivative of this mean
consists of terms with three overlap defects, and
\cref{lem:cav-threeproducts} bounds them by the third moment
$\mathcal M_3$. At the full endpoint $u=1$, exchangeability and the
contact identity $\E S_p=p$ show that the value of this mean is
$\Var(S_p)$, up to an error $1/(4n^2)$. Finally, \cref{lem:c7-thermal}
bounds $\mathcal M_3$ by $A$, $J$, $H_3$, and $N^{-3}$, and
\cref{cor:wi-cub} bounds $H_3$ by a small multiple of $V$ plus a
thermal term.

The rest of this section carries out this plan. In
\cref{sec:c7-coefficient-one}, we prove \cref{lem:cl-star} and the
trace identity of \cref{lem:cav-trace}. In \cref{sec:c7-third-moment},
we bound $\mathcal M_3$ in \cref{lem:c7-thermal}. In
\cref{sec:c7-closing}, we prove \cref{prop:c7-closure} and
\cref{cor:c7-defect}.

\section{Comparison with coefficient one and the trace identity}
\label{sec:c7-coefficient-one}

\begin{lemma}[Four-descendant variance comparison]\label{lem:cl-star}
Let $\mu=s\delta_0+\nu$ minimize $\mathcal F_{t,s}$. Then
\begin{equation}\label{eq:cl-star}
 V^-\le V+C_\beta\sup_{0\le u\le1}\mathcal M_3(u)+\frac1{4n^2}.
\end{equation}
\end{lemma}

\begin{proof}
For each level $p$, we represent $\E_u(A^-_p-p)^2$ as the mean of a
product of two overlap defects on four descendants. The proof has two
steps. In the first step, we compare the values of this mean at $u=0$
and $u=1$ using \cref{lem:streaming,lem:cav-threeproducts}. In the
second step, we evaluate it at $u=1$ using exchangeability and the
contact identity. Only the second step uses that $\mu$ is a minimizer.

\paragraph{Step 1: The interpolation.} Fix $p\in[0,1]$ and consider
four descendants that split at $p$. Each of the pairs $(1,2)$ and
$(3,4)$ splits at $p$, and we have $\Delta_{12}=R^-_{12}-p$ and
$\Delta_{34}=R^-_{34}-p$. We set
\[
 F_p=\Delta_{12}\Delta_{34}=(R^-_{12}-p)(R^-_{34}-p).
\]
Given $\calG_p$, the four spins are conditionally independent, each
with the conditional law of $\sigma$ given $\calG_p$ under $\Pp_u$. In
particular, $\E_{4,p,u}[R^-_{12}\mid\calG_p]=|M_p|^2/n=A^-_p$, and the
same holds for $R^-_{34}$. As the two pairs are conditionally
independent given $\calG_p$, we have
$\E_{4,p,u}[F_p\mid\calG_p]=(A^-_p-p)^2$. We define
\[
 \mathcal V_p(u)=\E_{4,p,u}F_p=\E_u(A^-_p-p)^2,\qquad
 \mathcal V(u)=\int\mathcal V_p(u)\,\nu(\dd p).
\]

Since $|F_p|\le4$, \cref{lem:streaming} implies that
$\mathcal V_p\in C^2([0,1])$, with $\mathcal V_p'(u)$ given by
\eqref{eq:streaming}. The right side of \eqref{eq:streaming} is
$\alpha/2$ times a signed combination of terms of two kinds. A term of
the first kind is the expectation, under a tree that extends the four
descendants, of $F_p$ times one further overlap defect and a spin factor
$\epsilon_{\mathsf x}\epsilon_{\mathsf y}$ of modulus one. The term of
the second kind is $s\,\mathcal V_p(u)\int\E_{2,r,u}D_{12}\,\nu(\dd r)$,
the product of $\mathcal V_p(u)$ with an expectation over an independent
pair of descendants. Each term contains three overlap defects, two of
which come from $F_p$. The absolute values of the coefficients, together
with $\nu(\dd p)$, form a mark measure of the kind considered in
\cref{subsec:c7-three}. Using \cref{lem:cav-threeproducts}, with four
initial leaves and one differentiation, and $\alpha\le\beta^2$, we find
\begin{equation}\label{eq:cl-star-derivative}
 \int|\mathcal V_p'(u)|\,\nu(\dd p)\le C_\beta\mathcal M_3(u)
 \qquad(0\le u\le1).
\end{equation}

For each $p$, we have
$\mathcal V_p(0)=\mathcal V_p(1)-\int_0^1\mathcal V_p'(u)\dd u$. At
fixed $N$, the derivative $\mathcal V_p'(u)$ is bounded uniformly in $p$
and $u$, because every term in \eqref{eq:streaming} is. We integrate
the identity against $\nu(\dd p)$, bound
$|\int_0^1\mathcal V_p'(u)\dd u|$ by $\int_0^1|\mathcal V_p'(u)|\dd u$,
and exchange the integrals in $p$ and $u$ by Tonelli's theorem. Using
\eqref{eq:cl-star-derivative}, we conclude that
\begin{equation}\label{eq:cl-star-step1}
 \mathcal V(0)\le\mathcal V(1)+C_\beta\sup_{0\le u\le1}\mathcal M_3(u).
\end{equation}
At $u=0$, the variance of $A^-_p$ under the rest law is at most its
second moment about $p$, which gives $V^-\le\mathcal V(0)$. Note that
this bound does not require $\E^-A^-_p=p$. We do not know whether this
identity holds, because the profile $\mu$ need not minimize the functional of
the rest system.

\paragraph{Step 2: The full endpoint.} We write
$f^{\rm site}_p=m_N(p)^2$. By \eqref{eq:c7-site-split}, we have
$nA^-_p=NS_p-f^{\rm site}_p$. Let $p\in\supp\nu$. By \eqref{eq:c7-exchange}
and the contact identity $\E S_p=p$ of \eqref{eq:in-contact}, we have
$\E_1f^{\rm site}_p=p$, and hence $\E_1A^-_p=(Np-p)/n=p$. Using
\eqref{eq:c7-exchange} again, we find
\[
\begin{aligned}
 \Var_1(A^-_p)
 &=\frac{N^2\Var(S_p)-2N\Cov_1(S_p,f^{\rm site}_p)
         +\Var_1(f^{\rm site}_p)}{n^2}\\
 &=\frac{N^2-2N}{n^2}\Var(S_p)+\frac{\Var_1(f^{\rm site}_p)}{n^2}
 \le\Var(S_p)+\frac1{4n^2}.
\end{aligned}
\]
In the last step, we used $N^2-2N=n^2-1$ and
$\Var_1(f^{\rm site}_p)\le1/4$, which holds because
$0\le f^{\rm site}_p\le1$. As $\E_1A^-_p=p$, the left side equals
$\mathcal V_p(1)$. Integrating against $\nu$, which is carried by
$\supp\nu$ and has mass at most one, we find
$\mathcal V(1)\le V+1/(4n^2)$. Together with \eqref{eq:cl-star-step1}
and $V^-\le\mathcal V(0)$, this proves \eqref{eq:cl-star}.
\end{proof}

The next lemma expresses $V$ itself through the cavity interpolation.
It uses exchangeability but not the minimizing property.

\begin{lemma}[Trace identity for the full variance]\label{lem:cav-trace}
Let the profile be deterministic with $\nu(\{1\})=0$, and let
\[
 V^{\rm site}=\int\Var_1\bigl(m_N(p)^2\bigr)\,\nu(\dd p).
\]
Then $0\le V^{\rm site}\le1/4$ and
\begin{equation}\label{eq:finaltrace}
 V=\frac{V^{\rm site}}N+\frac nN\Bigl[
 \frac\alpha d\Tr_\nu(J^{\scal}_{a,d}\Sigma^-)
 -\frac\alpha2\mathcal P+\alpha\int T_p\,\nu(\dd p)
 +\mathcal R_{\rm cav}\Bigr].
\end{equation}
\end{lemma}

\begin{proof}
Each variance in the definition of $V^{\rm site}$ lies in $[0,1/4]$,
because $0\le m_N(p)^2\le1$, and $\nu([0,1])\le1$; this proves the bounds
on $V^{\rm site}$.

Fix $p\in[0,1]$. By \eqref{eq:c7-exchange} and \eqref{eq:c7-site-split},
we have
\begin{equation}\label{eq:c7-trace-site}
 \Var(S_p)=\Cov_1\bigl(S_p,m_N(p)^2\bigr)
 =\frac1N\Var_1\bigl(m_N(p)^2\bigr)+\frac nN\,\mathcal C_{pp}(1).
\end{equation}
At $u=0$, the law $\Pp_0$ is the product of the law of the rest system
and the scalar law. The variable $m_N(p)$ depends only on the scalar
component and $A^-_p$ only on the rest component; hence
$\mathcal C_{pp}(0)=0$. Recall from \cref{subsec:c7-three} that
$\mathcal C_{pp}\in C^2([0,1])$. By Taylor's formula with integral
remainder, we have
\[
 \mathcal C_{pp}(1)=\mathcal C_{pp}'(0)
 +\int_0^1(1-u)\,\mathcal C_{pp}''(u)\dd u .
\]
At fixed $N$, all terms in this identity are bounded uniformly in $p$,
and we may integrate it against $\nu(\dd p)$. The integral of the
remainder is $\mathcal R_{\rm cav}$ by definition. As $\nu(\{1\})=0$,
\cref{lem:linear} applies, and \eqref{eq:lineartrace} gives the integral
of $\mathcal C_{pp}'(0)$. Integrating \eqref{eq:c7-trace-site} against
$\nu(\dd p)$, recalling that $V=\int\Var(S_p)\,\nu(\dd p)$, and inserting
these two integrals, we obtain \eqref{eq:finaltrace}.
\end{proof}

\section{The third moment of the rest overlap}\label{sec:c7-third-moment}

The remaining input to \cref{lem:remainder} and \cref{lem:cl-star} is
the integrated third moment $\mathcal M_3$ of the rest overlap defect.
We bound it by the thermal quantities $A$ and $J$ and the cubic moment
$H_3$ of the full system.

\begin{lemma}[Third moment of the rest overlap]\label{lem:c7-thermal}
For every deterministic profile,
\[
 \sup_{0\le u\le1}\mathcal M_3(u)
 \le K_{\beta,2}\bigl[64(A+J)+16H_3+32N^{-3}\bigr],
\]
where $K_{\beta,2}=e^{21\beta^2}$ is the constant of
\cref{cor:comparison} for two leaves, and $A$, $J$, and $H_3$ are
computed for the full system with the same profile.
\end{lemma}

\begin{proof}
The proof has three steps: an exact formula for the conditional second
moment of $R_q-S_q$ in the full system, the passage from the full
overlap to the rest overlap at $u=1$, and the comparison along the
interpolation.

\paragraph{Step 1: The full overlap.} Fix $q\in[0,1]$, and let
$\sigma^1,\sigma^2$ be conditionally independent given $\calG_q$, each
with the conditional law of $\sigma$, so that
$R_q=\sigma^1\cdot\sigma^2/N$. With $\delta^{\mathsf a}=\sigma^{\mathsf
a}-m_q$, we have
\[
 R_q-S_q=\frac{\delta^1\cdot\delta^2+m_q\cdot\delta^1+m_q\cdot\delta^2}N.
\]
Given $\calG_q$, the vectors $\delta^1$ and $\delta^2$ are independent
and centered, with covariance $C_q$. The three summands are pairwise
orthogonal in conditional $L^2$; for instance,
$\E[(\delta^1\cdot\delta^2)(m_q\cdot\delta^1)\mid\calG_q]=0$, because
$\delta^2$ is centered and independent of $\delta^1$. We find
\[
 \E\bigl[(R_q-S_q)^2\bigm|\calG_q\bigr]
 =\frac{\Tr C_q^2+2m_q^TC_qm_q}{N^2}.
\]
This is the conditional orthogonality in the proof of
\cref{lem:fnd-decomp} (see also \cite[Section~2.5, before (2.17)]{AL});
it does not use the contact identity and holds for every deterministic
profile. Taking expectations, we have
$\E(R_q-S_q)^2=a^{\rm th}_q+2j_q$. As $|R_q|\le1$ and $0\le S_q\le1$, we
have $|R_q-S_q|\le2$ and hence $|R_q-S_q|^3\le2(R_q-S_q)^2$. Integrating
against $\nu$, we conclude that
\begin{equation}\label{eq:c7-thermal-full}
 \int\E|R_q-S_q|^3\,\nu(\dd q)\le2(A+2J).
\end{equation}

\paragraph{Step 2: The rest overlap at $u=1$.} Since $\Pp_1$ is the full
tilted law, two descendants that split at $q$ have, under
$\E_{2,q,1}$, the joint law of the pair $(\sigma^1,\sigma^2)$ of
Step~1. Their full and rest overlaps satisfy
$NR_q=nR^-_{12}+\epsilon_1\epsilon_2$, which implies that
\[
 |R^-_{12}-R_q|=\frac{|R^-_{12}-\epsilon_1\epsilon_2|}N\le\frac2N .
\]
By the convexity of $x\mapsto|x|^3$, we have
$|x+y|^3\le4(|x|^3+|y|^3)$. Applying this first to
$(R^-_{12}-R_q)+(R_q-q)$ and then to $(R_q-S_q)+(S_q-q)$, we find
\[
 |R^-_{12}-q|^3\le32N^{-3}+16|R_q-S_q|^3+16|S_q-q|^3 .
\]
We take expectations and integrate against $\nu$, whose mass is at most
one. Using \eqref{eq:c7-thermal-full} and
$H_3=\int\E|S_q-q|^3\,\nu(\dd q)$, we obtain
\[
 \mathcal M_3(1)\le32N^{-3}+32A+64J+16H_3
 \le64(A+J)+16H_3+32N^{-3}.
\]

\paragraph{Step 3: Comparison.} For each $q$, the function
$|R^-_{12}-q|^3$ is nonnegative, bounded by $8$, and depends only on the
spins of two leaves that split at $q$. By \cref{cor:comparison} with
$k=2$, we have
$\E_{2,q,u}|R^-_{12}-q|^3\le K_{\beta,2}\E_{2,q,1}|R^-_{12}-q|^3$ for
every $u\in[0,1]$. Integrating in $q$ against $\nu$, we obtain
$\mathcal M_3(u)\le K_{\beta,2}\mathcal M_3(1)$, and the lemma follows
from Step~2.
\end{proof}

\section{Closing the variance inequality}\label{sec:c7-closing}

We now combine the trace identity, the error bounds, and the comparison
with coefficient one. The stability assumption below concerns the
scalar response operator of \cref{sec:sc} with the cavity coefficient
$d$ of \eqref{eq:cavityparameters}, acting on the scalar transport space
$\mathcal H_\nu$ of the minimizing profile.

\begin{proposition}[Self-overlap contraction]\label{prop:c7-closure}
Let $N\ge2$, $0<\rho\le1$, and $0<s\le1/2$, and let $\mu=s\delta_0+\nu$
minimize $\mathcal F_{t,s}$, with cumulative function $a$. Suppose that
the scalar response with coefficient $d=\beta^2(1-t/N)$ satisfies
\begin{equation}\label{eq:c7-stab}
 J^{\scal}_{a,d}\preceq(1+\eps_N)I\ \text{ on }\mathcal H_\nu,
 \qquad 0\le\eps_N\le\rho/4 .
\end{equation}
Then, with a constant $C_{\rm cav}$ depending only on $\beta$,
\begin{equation}\label{eq:c7-C1}
 \frac{3\rho}{4}\,V\le C_{\rm cav}\bigl[N^{-1}+A+J+H_3\bigr].
\end{equation}
\end{proposition}

\begin{proof}
The proof has three steps. We first collect the bounds on the error
terms in \eqref{eq:finaltrace} and on $V^-$. We then bound the trace
pairing using \eqref{eq:c7-stab}. Finally, we insert the comparison
\eqref{eq:cl-star} into \eqref{eq:finaltrace} and absorb the resulting
multiple of $V$.

\paragraph{Step 1: The errors.} As $\mu$ is a minimizer,
\eqref{eq:in-contact} implies that $\supp\nu\subset[0,1)$. Hence
$\nu(\{1\})=0$, and \cref{lem:linear,lem:cl-response,lem:cav-trace}
apply. Let
\[
 \mathfrak e=H_3+A+J+N^{-3}.
\]
We claim that
\begin{equation}\label{eq:c7-errors}
 \alpha\Bigl|\int T_p\,\nu(\dd p)\Bigr|\le C_\beta\bigl[A+J+N^{-2}\bigr],
 \qquad
 |\mathcal R_{\rm cav}|\le C_\beta\mathfrak e,
 \qquad
 V^-\le V+C_\beta\mathfrak e+\frac1{4n^2}.
\end{equation}
The first bound is \eqref{eq:cl-response-final}, with
$C_\beta=192\beta^2K_{\beta,3}$. For the second, by \cref{lem:remainder}, we have
$|\mathcal R_{\rm cav}|\le C_\beta\sup_w\mathcal M_3(w)$, and
\cref{lem:c7-thermal} bounds this supremum by $64K_{\beta,2}\mathfrak e$.
For the third, we insert the same bound on $\mathcal M_3$ into
\eqref{eq:cl-star}.

\paragraph{Step 2: The trace pairing.} Let $\bar A^-$ be the random
function $p\mapsto\bar A^-_p$. By \cref{lem:linear}, the trace pairing is
$\Tr_\nu(J^{\scal}_{a,d}\Sigma^-)=\E^-\langle\bar A^-,J^{\scal}_{a,d}\bar
A^-\rangle_{L^2(\nu)}$. The function $\bar A^-$ lies in
$\mathcal H_\nu$ for every realization, since $\bar A^-_0=0$ by
\eqref{eq:c7-root-overlap} and $\nu(\{1\})=0$. By Fubini's theorem, we
also have
$\E^-\|\bar A^-\|^2_{L^2(\nu)}=\int\Var^-(A^-_p)\,\nu(\dd p)=V^-$. By \cref{lem:Jdom}, $J^{\scal}_{a,d}$ is a
bounded self-adjoint operator on $\mathcal H_\nu$, and
\eqref{eq:c7-stab} states that
$\langle v,J^{\scal}_{a,d}v\rangle_{L^2(\nu)}\le(1+\eps_N)\|v\|^2_{L^2(\nu)}$
for every $v\in\mathcal H_\nu$. Taking $v=\bar A^-$ and averaging over
the rest system, we obtain
\[
 \Tr_\nu(J^{\scal}_{a,d}\Sigma^-)\le(1+\eps_N)V^-.
\]
Equivalently, we may test \eqref{eq:c7-stab} on the eigenvectors of
$\Sigma^-$ with nonzero eigenvalues, which lie in $\mathcal H_\nu$. We multiply the last display
by $\alpha/d$, which lies in $[0,t]$ by \eqref{eq:c7-alpha-d}. Using
$(1+\eps_N)V^-\ge0$, we find
\begin{equation}\label{eq:c7-trace-bound}
 \frac\alpha d\Tr_\nu(J^{\scal}_{a,d}\Sigma^-)\le t(1+\eps_N)V^-,
\end{equation}
whatever the sign of the trace.

\paragraph{Step 3: Absorption.} The term $-\alpha\mathcal P/2$ in
\eqref{eq:finaltrace} is nonpositive, because $\mathcal P\ge0$ by
\cref{lem:linear}. By \eqref{eq:c7-errors} and
\eqref{eq:c7-trace-bound}, the bracket in \eqref{eq:finaltrace} is at
most
\[
 t(1+\eps_N)V^-+C_\beta\bigl[A+J+N^{-2}+\mathfrak e\bigr],
\]
and this upper bound is nonnegative. As $0<n/N\le1$, the product of the
bracket with $n/N$ is also at most this upper bound. Using $V^{\rm site}\le1/4$ in \eqref{eq:finaltrace},
we have
\[
 V\le\frac1{4N}+t(1+\eps_N)V^-+C_\beta\bigl[A+J+N^{-2}+\mathfrak e\bigr].
\]
The stability assumption \eqref{eq:c7-stab} implies that
\[
 t(1+\eps_N)\le(1-\rho)\Bigl(1+\frac\rho4\Bigr)
 =1-\frac{3\rho}4-\frac{\rho^2}4\le1-\frac{3\rho}4\le1 .
\]
Inserting the bound on $V^-$ from \eqref{eq:c7-errors}, using
$t(1+\eps_N)V\le(1-3\rho/4)V$, and bounding $t(1+\eps_N)$ by one in the
error terms, we obtain
\[
 V\le\Bigl(1-\frac{3\rho}4\Bigr)V+\frac1{4N}+\frac1{4n^2}
 +C_\beta\bigl[A+J+N^{-2}+\mathfrak e\bigr].
\]
We move $(1-3\rho/4)V$ to the left side and note that each of
$(4n^2)^{-1}$, $N^{-2}$, and $N^{-3}$ is at most $N^{-1}$; for the first,
we used $n\ge N/2$. This proves \eqref{eq:c7-C1}, with a constant $C_{\rm cav}$
depending only on $\beta$.
\end{proof}

By \cref{cor:wi-cub}, $H_3$ is at most an arbitrarily small multiple of
$V$ plus a thermal term. Absorbing this multiple in \eqref{eq:c7-C1}, we
obtain the bound on the defect that we use in \cref{sec:b7}.

\begin{corollary}[Defect on the stable interval]\label{cor:c7-defect}
Let $N\ge2$, $1/N\le\rho\le1$, and $0<s=\lambda/\beta\le1/2$.
\begin{enumerate}[label=\textup{(\alph*)}]
\item Every minimizer $\mu=s\delta_0+\nu$ of $\mathcal F_{t,s}$ at these
parameters that satisfies \eqref{eq:c7-stab} for some $\eps_N$, with $a$
its cumulative function, obeys
\begin{equation}\label{eq:c7-defect}
 \beta(A+2J+V)\le C_\beta\frac{\log^2(2+N)}{N\lambda\rho^2},
\end{equation}
where $A$, $J$, and $V$ are computed at $\mu$, and $C_\beta$ depends only
on $\beta$.
\item If the defect $\calD_N(\rho)$ is defined at $\rho$ by
\eqref{eq:in-envelope} and some minimizer at $\rho$ satisfies
\eqref{eq:c7-stab}, then $\calD_N(\rho)$ satisfies the bound
\eqref{eq:c7-defect}.
\item If $1/N\le\rho_0\le1$ and, for every $\rho\in[\rho_0,1]$, every
minimizer at $\rho$ satisfies \eqref{eq:c7-stab}, then
\begin{equation}\label{eq:c7-integrated}
 \int_{\rho_0}^1\calD_N(\rho)\dd\rho
 \le C_\beta\frac{\log^2(2+N)}{N\lambda\rho_0}.
\end{equation}
\end{enumerate}
\end{corollary}

\begin{proof}
We first prove (a). Fix $C_{\rm cav}$ from \eqref{eq:c7-C1}, and let
$\kappa_\beta=\min\{1/4,(4C_{\rm cav})^{-1}\}$. As $\beta>1$ and
$\rho\ge1/N$, the parameter $M=Nc=\beta^2N\rho$ satisfies
$1<\beta^2\le M\le\beta^2N$, and $\kappa_\beta\rho\in(0,1/4]$. Hence the
hypotheses of \cref{cor:wi-cub} hold with $z=\kappa_\beta\rho$, and
\eqref{eq:wi-cub} gives
\[
 H_3\le\kappa_\beta\rho V
   +\frac{425\log(2M)}{Ms}\log\frac2{\kappa_\beta\rho}.
\]
We insert this bound into \eqref{eq:c7-C1}. The resulting term
$C_{\rm cav}\kappa_\beta\rho V$ is at most $\rho V/4$, and moving it to
the left side leaves $(\rho/2)V$ there. We find
\[
 \rho V\le2C_{\rm cav}\Bigl[N^{-1}+A+J
   +\frac{425\log(2M)}{Ms}\log\frac2{\kappa_\beta\rho}\Bigr].
\]
By \eqref{eq:in-Lvar} and \eqref{eq:th-J}, we have $A\le(Ms)^{-1}$ and
$J\le341(1+\log_+M)^2/(Ms)$. As $1\le M\le\beta^2N$ and $\rho\ge1/N$, we
also have
\[
 \log(2M)\le\log(2\beta^2)+\log N,\qquad
 1+\log_+M\le1+\log\beta^2+\log N,
\]
and $\log(2/(\kappa_\beta\rho))\le\log(2/\kappa_\beta)+\log N$. Each of
these three quantities is at most $C_\beta\log(2+N)$, because
$\log(2+N)\ge\log4>1$. Recalling that
$Ms=\beta^2N\rho\cdot\lambda/\beta=\beta N\rho\lambda$, we find
\[
 \beta V\le\frac{C_\beta}{N\rho}
   +\frac{C_\beta\log^2(2+N)}{N\lambda\rho^2},\qquad
 \beta(A+2J)\le\frac{C_\beta\log^2(2+N)}{N\lambda\rho}.
\]
Finally, $\rho\le1$, $\lambda=\beta s\le\beta/2$, and $\log(2+N)\ge1$
give
\[
 \frac1{N\rho}=\frac{\lambda\rho}{N\lambda\rho^2}
 \le\frac\beta2\cdot\frac{\log^2(2+N)}{N\lambda\rho^2},
 \qquad
 \frac{\log^2(2+N)}{N\lambda\rho}\le\frac{\log^2(2+N)}{N\lambda\rho^2},
\]
which proves \eqref{eq:c7-defect}.

For (b), recall from \eqref{eq:in-envelope} that, at every $\rho$ where
$\calD_N(\rho)$ is defined, $\calD_N(\rho)=\beta(A+2J+V)$ for every
minimizer at $\rho$. We evaluate the right side at a minimizer that
satisfies \eqref{eq:c7-stab}, and (a) gives the bound for
$\calD_N(\rho)$.

For (c), recall from \cref{sec:fnd} that $\calD_N$ is a nonnegative
measurable function defined at almost every level, and that a minimizer
exists at every $\rho$. Under the hypothesis of (c), part (b) applies at
almost every $\rho\in[\rho_0,1]$, and gives
$\calD_N(\rho)\le C_\beta\log^2(2+N)/(N\lambda\rho^2)$ there. This bound
holds at almost every level whichever minimizer is used, so no
measurable choice of minimizers is needed to integrate it. As
$\int_{\rho_0}^1\rho^{-2}\dd\rho=\rho_0^{-1}-1\le\rho_0^{-1}$,
integrating yields \eqref{eq:c7-integrated}.
\end{proof}
